\subsection{Классическая схема теории вероятностей}
Представьте, что вы играете в настольную игру и кидаете обычный игральный кубик. Каков шанс, что выпадет шестерка? Интуитивно понятно: один шанс из шести.
Это и есть "классическая схема теории вероятностей" (или классическое определение вероятности). Это самый простой, базовый и «честный» способ посчитать шанс наступления какого-то события. Если строго:
\begin{definition}[Классическая схема теории вероятностей]
Классическая схема теории вероятностей -- предположения, касаемые исследуемых экспериментов:
\begin{itemize}
    \item пространство элементарных исходов $\Omega$ -- конечное множество;
    \item элементарные исходы равновозможны.
\end{itemize}
\end{definition}

Последний пункт означает, что вероятность одинакова для каждого элементарного исхода. Тогда для событий:

\begin{definition}[Вероятность случайного события]
Вероятность случайного события $A$ в классической схеме -- отношение количества благоприятных исходов $N_A$ для события $A$ к общему числу исходов $N$:
\[
\mathbb{P}(A) = \frac{N_A}{N}.
\]
\end{definition}

Выделим довольно очевидные свойства:
\begin{enumerate}
    \item $\mathbb{P}(\varnothing) = 0$ и $\mathbb{P}(\Omega) = 1$;
    \item $\mathbb{P}(A \sqcup B) = \mathbb{P}(A) + \mathbb{P}(B)$;
    \item $\mathbb{P}(A \cup B) = \mathbb{P}(A) + \mathbb{P}(B) - \mathbb{P}(A \cap B)$;
    \item $\mathbb{P}(\overline{A}) = 1 - \mathbb{P}(\overline{B})$.
\end{enumerate}

Все задачи в классической схеме теории вероятностей сводятся так или иначе к комбинаторике, формулы которой будем считать известными.\\
\medskip
Кроме классического есть и другие определения вероятности. Их суть станет ясна немного позже:

\begin{definition}[Геометрическое определение вероятности]
Вероятность попадания случайной точки в заданную область $A$ равна отношению меры (длины, площади или объёма) этой области к мере всей возможной области исходов $\Omega$:
\[
\mathbb{P}(A) = \frac{\operatorname{mes}(A)}{\operatorname{mes}(\Omega)}.
\]
\end{definition}

\begin{definition}[Статистическое определение вероятности]
Вероятность события $A$ -- это предел, к которому стремится относительная частота наступления этого события при неограниченном увеличении числа одинаковых независимых испытаний:
\[
\mathbb{P}(A) = \lim_{n \to \infty} \frac{m}{n},
\]
где $m$ -- число появлений события, $n$ -- общее число испытаний.
\end{definition}

\begin{definition}[Аксиоматическое определение вероятности (по А.~Н.~Колмогорову)]
Вероятность -- это числовая функция (мера) $\mathbb{P}$, определённая на $\sigma$-алгебре событий пространства элементарных исходов $\Omega$, удовлетворяющая трём аксиомам:
\begin{enumerate}
    \item Неотрицательность: $\mathbb{P}(A) \ge 0$ для любого события $A$.
    \item Нормированность: $\mathbb{P}(\Omega) = 1$ (вероятность достоверного события равна 1).
    \item Счётная аддитивность: если события $A_1, A_2, \dots$ попарно несовместны, то
    \[
    \mathbb{P}\!\left(\bigcup_{i=1}^{\infty} A_i\right) = \sum_{i=1}^{\infty} \mathbb{P}(A_i).
    \]
\end{enumerate}
\end{definition}

Аксиоматическое определение является наиболее общим определением вероятности.